\chapter*{RESUMEN}
\thispagestyle{fancy}
\phantomsection
\addcontentsline{toc}{chapter}{RESUMEN}

Este trabajo presenta el diseño e implementación de un prototipo de robot móvil diferencial de bajo costo para mapeo 2D, navegación autónoma y exploración de interiores mediante \gls{ROS2} Foxy. El sistema integra una Raspberry Pi 4, un \gls{Arduino} Nano, encoders de cuadratura, un sensor RPLidar A1 y un gemelo digital desarrollado en Gazebo Classic. El control de velocidad utiliza un controlador PI cuyas ganancias fueron obtenidas mediante optimización por enjambre de partículas (PSO) y refinamiento experimental. El mapeo se implementó con SLAM Toolbox, la navegación con Nav2 y la exploración mediante una estrategia basada en fronteras.

La odometría presentó un error medio de 3.63\% en trayectorias rectas sobre hardware real. El control de velocidad fue evaluado a 0.13~m/s, obteniéndose un error estacionario cercano a cero y sin sobreimpulso apreciable dentro de la resolución de medida. La exploración autónoma produjo mapas válidos en cinco de las seis ejecuciones experimentales, con un índice medio de cobertura libre de 62.2\% y un área libre mapeada media de 3.50~m$^2$. Durante estas ejecuciones se generaron 60 metas de navegación, de las cuales 54 fueron abortadas, una fue completada y cinco quedaron sin resultado terminal. Los resultados verifican funcionalmente la integración del proceso de mapeo y exploración autónoma, aunque evidencian como principal limitación la baja robustez del sistema para completar individualmente las metas de navegación generadas.

\noindent\textbf{Palabras clave:} robot móvil diferencial, SLAM, \gls{ROS2}, Nav2, navegación autónoma, exploración de fronteras, PSO, gemelo digital.

\chapter*{ABSTRACT}
\thispagestyle{fancy}
\phantomsection
\addcontentsline{toc}{chapter}{ABSTRACT}

This thesis presents the design and implementation of a low-cost differential mobile robot prototype for 2D mapping, autonomous navigation, and indoor exploration using \gls{ROS2} Foxy. The system integrates a Raspberry Pi 4, an \gls{Arduino} Nano, quadrature encoders, an RPLidar A1 sensor, and a digital twin developed in Gazebo Classic. Speed control uses a PI controller whose gains were obtained through Particle Swarm Optimization (PSO) and experimental refinement. Mapping was implemented with SLAM Toolbox, navigation with Nav2, and exploration through a frontier-based strategy.

Odometry achieved a mean error of 3.63\% in straight-line trajectories on real hardware. Speed control was evaluated at 0.13~m/s, showing near-zero steady-state error and no appreciable overshoot within the measurement resolution. Autonomous exploration produced valid maps in five of the six experimental runs, with a mean free-space coverage index of 62.2\% and a mean mapped free area of 3.50~m$^2$. During these runs, 60 navigation goals were generated, of which 54 were aborted, one was completed, and five remained without a terminal result. The results functionally verify the integration of the autonomous mapping and exploration process, while identifying the limited robustness of the system in completing individual navigation goals as the main performance limitation.

\noindent\textbf{Keywords:} differential mobile robot, SLAM, \gls{ROS2}, Nav2, autonomous navigation, frontier exploration, PSO, digital twin.
